\section{Worked compilation examples}\label{sec:examples}

\subsection{Example A: a residual height-bounded enumerator}

Let \(P(c,\vec x)\) be a reflected predicate on canonical point codes, let \(h(c)\) be a primitive-recursive height, and let \(H(\vec x)\) be a proposed cutoff.  Define \(L_H\) by enumerating all codes in a fixed box determined by \(H\), filtering by well-formedness and \(P\), canonicalizing, and deduplicating.

The graph has nodes for box generation, predicate reflection, filtering, canonicalization, deduplication, soundness, and completeness.  Soundness is unconditional:
\[
  c\in L_H\longrightarrow P(c,\vec x)=1.
\]
Completeness consumes the boundary formula
\[
 \mathsf{HeightBound}(\vec x,H):
 \quad P(c,\vec x)=1\longrightarrow h(c)\le H(\vec x).
\]
The query for soundness alone reaches no height boundary.  The query for completeness reaches exactly \(\mathsf{HeightBound}\).  The combined query shares enumeration and filtering.  The compiler exports
\[
 \Contract([\mathsf{HeightBound}],\mathsf{Sound}(L_H)\wedge
          \mathsf{Complete}(L_H)).
\]
This demonstrates query-minimal residuals: a verifier does not prove the existence of the cutoff it checks.

\subsection{Example B: exact rational bisection of \(y^3+y=x\)}

For rational \(x\) and an initial interval \([a,b]\) with
\[
 a^3+a\le x\le b^3+b,
\]
run \(N\) bisection steps.  The state stores the rational endpoints.  At each step, compare the midpoint polynomial with \(x\), retaining the half interval containing a root.  Monotonicity follows from
\[
 (v^3+v)-(u^3+u)=(v-u)(u^2+uv+v^2+1),
\]
and the second factor is positive for ordered rationals.

All comparisons are exact integer inequalities after clearing positive denominators.  The loop invariant is the conjunction of endpoint order, bracketing, and the denominator convention; after Booleanization it is quantifier-free.  The normalized trace kernel proves the invariant at the computed bound.  A final arithmetic contract proves interval width \((b-a)/2^N\).  There is no branch-separation residual because equality at the midpoint may be handled by either closed half interval.  The public output is an interval, not a discontinuously selected sign branch.

\subsection{Example C: a bounded challenge--response theorem}

Suppose \(A(u,v,w,\vec x)\) is reflected and all bounds are primitive recursive.  Consider
\[
 F(\vec x)=\forall u<U\;\exists v<V(u)\;\forall w<W(u,v)\;A(u,v,w,\vec x).
\]
The strategy compiler constructs a truth certificate consisting of a response table \(v(u)\) and atomic truth data for every legal challenge \((u,w)\).  A falsity certificate chooses one \(u\), gives for every candidate \(v\) a counterchallenge \(w(v)\), and supplies atomic falsity data.  The two checkers are reflected by structural contracts.  A later theorem can consume \(F\) without knowing which finite strategy representation was used.

\subsection{Example D: a dead-end relation}

Let levels be
\[
 X_0=\{0,1\},\quad X_1=\{0,1,2\},\quad X_2=\{2\},
\]
with edges \(0\to0\), \(1\to1\), \(1\to2\), and \(2\to2\) where types permit.  Backward viability at horizon 2 removes \(0\in X_1\), then removes \(0\in X_0\); only the branch beginning at \(1\) survives.  Forward greedy choice without viability could choose \(0\) and get stuck.  The reference implementation checks that the least viable path is \((1,2,2)\).

\subsection{Example E: cyclic precision requests}

Take two requests \(r_0,r_1\) with seeds \((1,0)\), caps \((8,8)\), and updates
\[
 r_0\ge r_1+1,
 \qquad
 r_1\ge \lfloor r_0/2\rfloor+1.
\]
The capped monotone iteration stabilizes at the least closed vector above the seeds.  The rank
\[
 (8-r_0)+(8-r_1)
\]
strictly decreases on every nonstationary update, so at most 17 rounds are needed.  If a transformer requests a value above its target cap, the overflow Boolean identifies the offending edge; the compiler turns cap adequacy into a residual rather than silently truncating a required precision.

\subsection{Example F: finite probability}

A finite rational law is represented by weights \(w_i\), a positive total \(W\), and values \(X_i\).  Expectations are carried division-free as
\[
                 E_X=\sum_i w_iX_i,
\]
with the normalized value understood as \(E_X/W\) only under the guard \(W>0\).  The checked finite-sum and Fubini backend proves linearity, product-law identities, tail sums, and finite total expectation.  Approximation to an infinite law is a separate contract requiring an explicit tail-moment or metric error bound.  The finite theorem has no hidden measure-theoretic compactness.

\subsection{Example G: root clusters}

For a polynomial with approximate coefficients, a finite complex grid and argument-principle certificate can identify disjoint boxes with multiplicities.  The stable output is the multiset of boxes and multiplicities.  A query for labeled roots adds separation and continuation residuals; a query for symmetric polynomials does not.  This is precisely the benefit of query specialization: the stronger, discontinuous output does not burden the invariant output.
